\chapter{Mimari Arka Plan: SAM'den EdgeTAM'a}
\label{ch:arka-plan}

Bu bölüm, bu çalışmanın hızlandırdığı modelin nereden geldiğini ve
\emph{neden} bu şekilde tasarlandığını anlatır. Amaç SAM ailesinin genel
bir tanıtımı değil, EdgeTAM'ı hızlı kılan tasarım kararının --
memory attention'ın token sayısını küçültmek -- somut sayılarla
gösterilmesidir; çünkü bu proje, tam olarak o küçültülmüş mimariyi
TensorRT'e taşımaktadır.

\section{SAM: Statik Görüntüde Komutla Segmentasyon}

\emph{Segment Anything Model} (SAM)~\cite{kirillov2023sam}, tek bir
görüntü üzerinde bir komuta (bir kutu, bir nokta) karşılık gelen nesnenin
maskesini üreten üç parçalı bir mimaridir:

\begin{itemize}
  \item \textbf{image encoder}: görüntüyü bir kez işleyip yoğun bir
        özellik haritası üretir (bir Vision Transformer);
  \item \textbf{prompt encoder}: kullanıcının işaretlediği nokta/kutuyu
        bir embedding vektörüne çevirir;
  \item \textbf{mask decoder}: özellik haritasını ve prompt embedding'ini
        birleştirip maskeyi üretir -- hafif bir modüldür, çünkü ağır iş
        zaten image encoder'da yapılmıştır.
\end{itemize}

SAM'in tasarımı görüntü başına \emph{bir kez} çalışacak şekilde
optimize edilmiştir: image encoder pahalıdır ama tek seferliktir,
komut değiştikçe yalnızca ucuz decoder tekrar çalışır. Bu tasarım
video için doğrudan uygun değildir -- her kare için sıfırdan
çalıştırılırsa hem yavaştır hem de kareler arasında hiçbir bilgi
taşımaz; nesne bir karede kaybolsa (örn. kısa süreliğine engellenmiş
olsa) bir sonraki karede yeniden bulunamaz.

\section{SAM 2: Video ve Memory Mekanizması}
\label{sec:sam2-bellek}

\emph{SAM 2}~\cite{ravi2024sam2}, SAM'in imza mimarisine bir
\textbf{memory bank} ekleyerek video için genişletir. Fikir basittir:
model her karede ürettiği maskeyi ve o karenin görsel özelliklerini bir
``hafızaya'' yazar; sonraki karede bu hafızayı okuyarak nesnenin nerede
olması gerektiğine dair bir önceliği (prior) elde eder. Üç yeni parça
eklenir:

\begin{itemize}
  \item \textbf{memory bank}: son birkaç karenin (bu projede
        kullanılan kontrol noktasında en fazla 7 -- \texttt{num\_maskmem})
        özellik haritalarını ve object pointer'larını tutan, kare
        indeksiyle anahtarlanmış bir yapı;
  \item \textbf{memory attention}: geçerli karenin özelliklerini
        memory bank'takilerle karşılaştıran cross-attention katmanları;
  \item \textbf{memory encoder}: geçerli karenin maskesini ve
        özelliklerini memory bank'a yazılacak biçime dönüştürür.
\end{itemize}

\subsection{memory attention'ın gerçek maliyeti}

SAM 2'nin kontrol noktası 1024×1024 girdiyle çalışır; image encoder'ın
stride oranı sonucu her kare $H{=}W{=}64$, yani \textbf{4096 uzamsal
token} olarak temsil edilir. Sorun, memory attention'ın bu token'ları
\emph{nasıl} karşılaştırdığındadır: SAM 2 her hafıza karesini tam
çözünürlükte -- yani yine 4096 token olarak -- saklar. $T$ kare hafızada
tutulduğunda, geçerli karenin 4096 sorgu token'ı bu $T{\times}4096$
anahtar/değer token'ıyla karşılaştırılır. Dikkat mekanizmasının
maliyeti sorgu sayısı ile anahtar sayısının çarpımıyla orantılı
olduğundan, EdgeTAM makalesinin de belirttiği gibi karmaşıklık

\begin{equation}
  \mathcal{O}(T \, C \, H^2 W^2)
  \label{eq:sam2-karmasiklik}
\end{equation}

biçimindedir ($C$: kanal/gizli boyut). Buradaki $H^2W^2$ terimi kritik
öneme sahiptir: çözünürlüğü iki katına çıkarmak dikkat maliyetini dört
değil, \textbf{on altı} katına çıkarır -- hem sorgu hem anahtar tarafı
karesel olarak büyüdüğü için. Bu, SAM 2'nin kenar cihazlarda neden
saniyede tek kareye kadar yavaşladığının matematiksel kökenidir ve
EdgeTAM makalesinin de doğrudan tespit ettiği nokta: yalnızca image
encoder'ı küçültmek beklenen hızlanmayı getirmez, çünkü darboğaz orada
değildir.

\section{EdgeTAM: 2D Spatial Perceiver ile Memory Sıkıştırma}
\label{sec:edgetam-perceiver}

EdgeTAM'ın~\cite{zhou2025edgetam} temel katkısı, her hafıza karesini
\emph{tam çözünürlükte} saklamak yerine sabit sayıda öğrenilebilir
latent token'a sıkıştırmaktır. Bu iş bir \textbf{2D Spatial Perceiver}
tarafından yapılır ve iki tür token üretir:

\begin{itemize}
  \item $N_g{=}256$ \textbf{global latent token} (\emph{Global
        Perceiver}): her biri karenin tamamına bakıp genel bir özet
        çıkarır, uzamsal konum taşımaz;
  \item $N_l{=}256$ \textbf{2D spatial latent token}: karenin uzamsal
        düzenini (hangi bölgede ne olduğunu) daha kaba bir ızgarada
        korur.
\end{itemize}

Toplamda hafızaya yazılan her kare $N_g{+}N_l{=}\mathbf{512}$ token'a
indirgenir -- 4096'ya karşı \textbf{8 kat} azalma. memory attention'ın
anahtar/değer tarafı artık $T{\times}4096$ değil $T{\times}512$
büyüklüğündedir, dolayısıyla karmaşıklık

\begin{equation}
  \mathcal{O}(T \, C \, H W \, (N_g{+}N_l))
  \label{eq:edgetam-karmasiklik}
\end{equation}

biçimine düşer: $H^2W^2$'deki iki $HW$ çarpanından biri sabit bir sayıyla
($N_g{+}N_l$) yer değiştirmiştir. EdgeTAM makalesi bu değişikliğin
memory attention modülünü tek başına \textbf{8×} hızlandırdığını
raporlar; ayrıca dikkat katman sayısı SAM 2'nin 4'ünden \textbf{2}'ye
indirilmiştir. İki değişikliğin bir aradaki etkisiyle EdgeTAM, iPhone 15
Pro Max'te SAM 2'ye göre \textbf{22×} hızlanma (0,7~FPS → 16~FPS)
bildirir.

\begin{figure}[H]
  \centering
  \begin{tikzpicture}[font=\sffamily\small]
    \def\barw{0.9}
    % SAM2 bar - very tall (4096)
    \draw[fill=grimetin!25, draw=grimetin] (0,0) rectangle (\barw, 4.096);
    \node at (\barw/2, 4.096+0.35) {\textbf{4096}};
    \node at (\barw/2, -0.5) {\shortstack{SAM 2\\(tam çözünürlük)}};
    % EdgeTAM bar - short (512), drawn to same scale
    \draw[fill=anarenk!30, draw=anarenk] (2.2,0) rectangle (2.2+\barw, 0.512);
    \node at (2.2+\barw/2, 0.512+0.35) {\textbf{512}};
    \node at (2.2+\barw/2, -0.5) {\shortstack{EdgeTAM\\($N_g{+}N_l$)}};
    % axis
    \draw[-{Latex[length=2mm]}] (-0.6,0) -- (-0.6,4.4) node[above, font=\sffamily\footnotesize]{token/kare};
    \node[anchor=west, color=vurgu, font=\sffamily\bfseries] at (3.6, 2.0) {8× azalma};
    \draw[-{Latex[length=1.5mm]}, vurgu] (3.5,1.85) -- (2.2+\barw+0.1, 1.0);
  \end{tikzpicture}
  \caption{Hafızaya yazılan kare başına token sayısı: SAM 2'nin tam
    uzamsal çözünürlüğüne karşı EdgeTAM'ın Perceiver ile sıkıştırılmış
    hafızası. memory attention'ın anahtar/değer tarafı bu sayıyla
    doğru orantılıdır.}
  \label{fig:token-karsilastirma}
\end{figure}

\subsection{Knowledge Distillation}

Sıkıştırılmış hafıza, SAM 2'nin sıfırdan eğitilmiş bir küçük modeliyle
değil, \textbf{knowledge distillation} ile elde edilir: SAM 2 öğretmen,
EdgeTAM öğrenci rolündedir. Eğitim iki aşamalıdır -- önce tekil
görüntüler üzerinde öğretmenle aynı segmentasyon çıktısını üretmeyi
öğrenme, ardından video dizileri üzerinde öğretmenin memory/takip
davranışını taklit etme. Bu, EdgeTAM'ın küçük bir omurga (RepViT
tabanlı) ve sıkıştırılmış bir hafıza ile SAM 2'ye yakın doğruluğa
ulaşmasını sağlayan eğitim tarafındaki tamamlayıcı unsurdur; bu raporun
konusu olan inference hızlandırması bu eğitimden sonraki, sabit
ağırlıklı kontrol noktası üzerinde çalışır.

\section{Bu Projenin EdgeTAM Mimarisiyle İlişkisi}
\label{sec:proje-eslesme}

Bu çalışma, yukarıda anlatılan mimariyi değiştirmez -- eğitilmiş
EdgeTAM kontrol noktasını alıp dört hesaplama bloğunu TensorRT motoruna
dönüştürür. \Cref{tab:flop-dagilimi-arkaplan}, bu dört bloğun gerçek
kontrol noktası üzerinde 1024×1024 çözünürlükte, tek nesneyle ölçülen
işlem yükü (FLOP) ve parametre sayılarını gösterir
(\texttt{tools/analyze\_cost.py}, ayrıntılar \cref{sec:flop-dagilimi}).

\begin{table}[H]
  \centering
  \caption{Kare başına hesaplama dağılımı, gerçek kontrol noktası,
    1024×1024. \emph{(CPU'da, yapısal -- FLOP donanımdan bağımsızdır.)}}
  \label{tab:flop-dagilimi-arkaplan}
  \begin{tabular}{lrrrl}
    \toprule
    Modül & GFLOP/kare & Pay & Parametre (M) & Bu bölümde nereye karşılık gelir \\
    \midrule
    image encoder & 38,7 & \%26,8 & 4,92 & \S\ref{sec:sam2-bellek} -- SAM'in encoder'ı \\
    \textbf{memory attention} & \textbf{89,3} & \textbf{\%61,9} & 2,96 & \S\ref{sec:edgetam-perceiver} -- sıkıştırılmış dikkat \\
    memory encoder + Perceiver & 12,6 & \%8,8 & 1,62 & \S\ref{sec:edgetam-perceiver} -- sıkıştırmanın kendisi \\
    SAM head & 3,6 & \%2,5 & 4,41 & \S 2.1 -- SAM'in mask decoder'ı \\
    \midrule
    \textbf{Toplam} & \textbf{144,4} & & 13,9 & \\
    \bottomrule
  \end{tabular}
\end{table}

Dikkat çekici nokta parametre sütunudur: memory attention ağırlıkların
yalnızca \%21'ini oluşturur ama işlem yükünün \%62'sidir -- model
boyutu, zamanın nereye gittiğinin güvenilir bir göstergesi değildir.
Bu tam olarak \cref{ch:giris}'te bahsedilen, ``yalnızca image encoder'ı
hızlandırmanın yetmediği'' bulgusunun kaynağıdır ve bu çalışmanın dört
motora -- yalnızca bire değil -- neden ihtiyaç duyduğunu doğrudan
açıklar (\cref{ch:sistem-yontem}).
